% Architecture layer ribbon snippets for the final presentation.
% Use the macro once before the detailed architecture layer slides.
% Then add \ArchitectureLayerRibbon{1}, ..., \ArchitectureLayerRibbon{5}
% near the top of each corresponding layer slide.

% ---------------------------------------------------------------------------
% Mini architecture ribbon for the detailed layer slides.
% The highlighted number is the current architecture layer.
\newcommand{\ArchitectureLayerRibbon}[1]{%
\begin{center}
\begin{tikzpicture}[
  font=\scriptsize,
  archbox/.style={
    draw,
    rounded corners=3pt,
    minimum width=2.45cm,
    minimum height=0.72cm,
    align=center,
    line width=0.8pt
  },
  arr/.style={-Stealth, thick, gray!55}
]

\node[archbox,
      fill=\ifnum#1=1 phase1blue!18\else gray!10\fi,
      text=\ifnum#1=1 phase1blue!80!black\else black\fi] (a1)
      at (0,0) {1. Research\\Context};

\node[archbox,
      fill=\ifnum#1=2 phase2green!18\else gray!10\fi,
      text=\ifnum#1=2 phase2green!60!black\else black\fi] (a2)
      at (2.9,0) {2. Control +\\Validation};

\node[archbox,
      fill=\ifnum#1=3 phase3orange!18\else gray!10\fi,
      text=\ifnum#1=3 phase3orange!80!black\else black\fi] (a3)
      at (5.8,0) {3. Execution};

\node[archbox,
      fill=\ifnum#1=4 phase4purple!18\else gray!10\fi,
      text=\ifnum#1=4 phase4purple!80!black\else black\fi] (a4)
      at (8.9,0) {4. Policy + Allocation\\+ Mediation};

\node[archbox,
      fill=\ifnum#1=5 phase4purple!10\else gray!10\fi,
      text=\ifnum#1=5 phase4purple!80!black\else black\fi] (a5)
      at (12.0,0) {5. Artifacts +\\Reporting};

\draw[arr] (a1) -- (a2);
\draw[arr] (a2) -- (a3);
\draw[arr] (a3) -- (a4);
\draw[arr] (a4) -- (a5);

\end{tikzpicture}
\end{center}
}
% ---------------------------------------------------------------------------

% Example insertion points:
%
% \begin{frame}{Architecture Layer 1: Research Context}
% \ArchitectureLayerRibbon{1}
% ... existing Layer 1 content ...
% \end{frame}
%
% \begin{frame}{Architecture Layer 2: Control and Validation}
% \ArchitectureLayerRibbon{2}
% ... existing Layer 2 content ...
% \end{frame}
%
% \begin{frame}{Architecture Layer 3: Execution}
% \ArchitectureLayerRibbon{3}
% ... existing Layer 3 content ...
% \end{frame}
%
% \begin{frame}{Architecture Layer 4: Policy, Allocation, and Mediation}
% \ArchitectureLayerRibbon{4}
% ... existing Layer 4 content ...
% \end{frame}
%
% \begin{frame}{Architecture Layer 5: Artifacts and Reporting}
% \ArchitectureLayerRibbon{5}
% ... existing Layer 5 content ...
% \end{frame}

% ---------------------------------------------------------------------------
% Optional one-slide version if there is only time for one layer explanation.
\begin{frame}{Architecture Layers as Domain Phases}
\small

\begin{block}{How to read the architecture}
Each detailed layer slide can show this mini-ribbon at the top, with the current layer highlighted. The five implementation layers still follow the same four-phase logic: environment, context, decision, and evaluation.
\end{block}

\vspace{0.2em}
\ArchitectureLayerRibbon{5}
\vspace{0.1em}

\begin{tabular}{p{0.20\textwidth}p{0.27\textwidth}p{0.43\textwidth}}
\toprule
\textbf{Domain phase} & \textbf{Architecture layer} & \textbf{What to point out} \\
\midrule
\rowcolor{phase1blue!12}
Phase 1: Environment & Layer 1 + Layer 3 & Research context defines the problem; execution turns clinical and quantum testbeds into runnable cases. \\

\rowcolor{phase2green!12}
Phase 2: Context & Layer 2 + Layer 3 & Control and validation preserve run settings, context quality, missingness, and provenance. \\

\rowcolor{phase3orange!12}
Phase 3: Decision & Layer 4 & Models, allocators, and EQUITAS mediation choose or adjust routing actions under partial feedback. \\

\rowcolor{phase4purple!12}
Phase 4: Evaluation & Layer 5 & Fairness states, summaries, manifests, and validation-hub artifacts turn runs into evidence. \\
\bottomrule
\end{tabular}

\vspace{0.3em}

\begin{block}{Fast speaking point}
If time is short, focus on \textbf{Layer 5}: it shows why the project is reproducible, because claims link back to states, artifacts, manifests, and validation evidence.
\end{block}

\slidefoot{Acronyms: EQUITAS = fairness audit and mitigation boundary. Sources: DSCI601 architecture report, approach report, workspace fairness workflow, artifact writer, and validation hub.}
\end{frame}
